% !PRINTED 532
\weberpage{532}
Es gilt aber auch in Bezug auf die \gesperrt{Functionale} der Satz:

2. \gesperrt{Wenn $\omega_1, \omega_2, \ldots, \omega_n$ eine Basis von $\mathfrak{o}$ ist, so sind in der Form}
\weq{9}{\omega = u_1\omega_1 + u_2\omega_2 + \cdots + u_n\omega_n,}
\gesperrt{in der $u_1, u_2, \ldots, u_n$ ganze rationale Functionale sind, alle ganzen Functionale in $\Omega$ enthalten.}

Denn wir haben schon oben gezeigt, dass alle Functionale überhaupt in der Form (9) enthalten sind, wenn die Coëfficienten $u_1, u_2, \ldots, u_n$ ganze oder gebrochene rationale Functionale sind. Wir können aber immer eine ganze primitive Function $e$ so bestimmen, dass
\[
e u_1 = y_1, \; e u_2 = y_2, \; \ldots, \; e u_n = y_n
\]
ganze Functionen der Variablen sind, und dann wird
\weq{10}{e\omega = y_1\omega_1 + y_2\omega_2 + \cdots + y_n\omega_n,}
und da $e$ eine Einheit ist, so ist $e\omega$ zugleich mit $\omega$ ganz. Da nun die Coëfficienten der Potenzen und Producte der Variablen in der Function (10) nach \S.\,142, 3. ganze Zahlen sein müssen, so folgt nach 1., dass die Coëfficienten in den Functionen $y_1, y_2, \ldots, y_n$ ganze rationale Zahlen sein müssen, und dass folglich $u_1, u_2, \ldots, u_n$ ganze rationale Functionale sind.

\begin{center}
\S.\,146.

\gesperrt{Die Basen der Functionale.}
\end{center}

\gesperrt{Ist $\mu$ ein ganzes Functional des Körpers $\Omega$, so verstehen wir unter einer Basis von $\mu$ ein System von $n$ ganzen Zahlen in $\Omega$:}
\weq{1}{\alpha_1, \alpha_2, \ldots, \alpha_n,}
\gesperrt{das eine Basis des Körpers $\Omega$ ist, und dem die Eigenschaft zukommt, dass in der Form}
\weq{2}{\alpha = x_1\alpha_1 + x_2\alpha_2 + \cdots + x_n\alpha_n}
\gesperrt{alle durch $\mu$ theilbaren ganzen Zahlen des Körpers $\Omega$ und keine anderen enthalten sind, wenn für $x_1, x_2, \ldots, x_n$ ganze rationale Zahlen gesetzt werden.}

Es soll jetzt bewiesen werden, dass jedes ganze Functional eine Basis hat.
